\documentclass[11pt]{article}

\usepackage{amsmath,amssymb,amsthm}
\usepackage{graphicx}
\usepackage{hyperref}
\usepackage{geometry}
\geometry{margin=1in}

\usepackage{macros}

\title{Variations of Apple-Peel Unfoldings}
\author{Takashi Yoshino}
\date{04/04/2026}

\begin{document}
\maketitle

\begin{abstract}
"Apple-peel unfolding" is the procedure to obtain unfolded expansions of polyhedra. This article describes the variation in its definitions and compares the unfolded patterns. 
\end{abstract}

\section{Introduction}
This is a sample introduction.

\section{Three Definitions and Algorithms}
Peeling apples is intuitive. Some tried to peel 




A loxodrome (also called a rhumb line) is a curve on a sphere that crosses
  every meridian (line of longitude) at the same angle.

  On a globe: if you sail from one point to another while keeping your compass
  bearing constant — say, always heading 45° northeast — the path you trace is
  a loxodrome. It spirals around toward the pole rather than taking the
  shortest path (which would be a great circle).

  Key properties:
  - On a Mercator map projection, a loxodrome appears as a straight line — this
   is exactly why Mercator projection was invented for navigation
  - It is NOT the shortest path between two points (that would be a great
  circle / geodesic)
  - It spirals infinitely around the poles without ever reaching them (except
  when heading due north/south)


 Why the name "R2 (loxodrome)" here:

  In this peeling algorithm, the face centroids lie approximately on a sphere.
  Rule R2 selects the candidate with the smallest angular deviation from the
  forward direction — meaning it tries to continue as straight as possible at
  each step. This mimics the behavior of a loxodrome: maintaining a nearly
  constant bearing while spiraling around the sphere.

  In contrast:
  - R1 (max φ) takes the sharpest left turn — like a tight spiral
  - R2 (loxodrome) takes the gentlest left turn — like a gradual rhumb-line
  spiral
  - R3 (max z) climbs toward the top as fast as possible — like following a
  meridian upward

  The name reflects the geometric intuition, not an exact correspondence.



\section{Main Results}

\begin{theorem}
Let $M$ be a compact manifold. Then something interesting happens.
\end{theorem}

\begin{proof}
This is a placeholder proof.
\end{proof}

\bibliographystyle{plain}
\bibliography{references}

\end{document}
